\section*{September 22, 2026: Quarter-mile results on the new sample: no average effect, a gradient in baseline price}
\textit{Drafted by Claude. Jacob chose the quarter-mile focus, log price as the outcome, and fixed effects plus hedonics as the main specification.}

On the September 22 clean sample, closures have no detectable average effect on
sale prices or sales volume within a quarter mile. The average hides a gradient.
Relative to stayed-open schools, homes near closed schools did better after 2013
the more expensive their neighborhood was before the closures. Part of that gradient
lines up with gentrification already under way before 2013, and part does not.

The comparison is the September 22 clean sample within 1,320 feet of the 29
closed-program sites and 48 stayed-open sites. It excludes sales near both groups, near welcoming
schools, near other February candidates, and near buildings vacated by
relocating welcoming schools. The outcome is log real price. The main
specification has site and year fixed effects plus hedonics: size, lot, age,
rooms, beds, baths, residence type, quality, condition, REO and quick-resale
indicators, and a two-to-six-unit indicator. It has no property-class dummies.
All clean sales enter, with equal weight per sale, and errors are clustered by site. The pooled
estimate compares 2014--2018 with 2008--2012 and omits 2013. It is 0.050
[$-0.12$, 0.23] for all property (11,042 sales) and $-0.030$ [$-0.16$, 0.10] for
single-family homes (5,384). Neither rejects a flat pre-period (pre-trend
$p = 0.90$ and 0.74). The estimates are similar after dropping REO resales, quick
resales, price tails, or all flagged sales. The largest is 0.112 [$-0.13$, 0.36], for all
property without REO resales. Equal site-year weights give $-0.010$ and $-0.076$.
A unit-count factor with an unknown level moved no estimate by more than 0.003.

The single-family 2010 dip seen without controls is not a pre-trend. Relative to
the 2008--2012 average rather than to 2012 alone, the hedonic single-family series
alternates between years (2010 is $+0.06$). Without controls, 2012--2014 are high
years, not 2010 a low one. Dropping each closed site in turn moves the hedonic
2010 coefficient only between $-0.04$ and 0.07.

Sales volume shows no effect either. Poisson site-year models count all market sales,
clean sales, clean sales without REO, and REO resales. They give pooled estimates between $-0.13$ and
0.01, none distinguishable from zero. Closed-school areas are thinner markets throughout;
there is no parcel-count denominator, so these are sales, not turnover rates.

\begin{center}
\includegraphics[page=1,width=0.9\textwidth]{../input/main_spec_2026-09-22.pdf}
\end{center}

Splitting sites at the median of their 2008--2012 median sale prices gives
opposite signs. The estimate is 0.153 [$-0.02$, 0.32] in higher-price areas (15 closed, 23
stayed-open sites) and $-0.120$ [$-0.38$, 0.14] in lower-price areas (14, 25).
Quartiles give $-0.06$, $-0.19$, 0.02 and 0.20 [$-0.00$, 0.41], for all property.
A continuous interaction with the log site median price also lets year effects
vary with baseline price for every site. It gives a slope of 0.159 [0.01, 0.31] per log
point for all property and 0.117 [$-0.00$, 0.24] for single-family homes. The
quartile estimates lie on that line. Year by year, the all-property slope is about zero
in 2008--2011 and $+0.06$ to $+0.17$ from 2013 on. Volume shows no comparable
split: in higher-price areas it is $-0.08$ [$-0.20$, 0.04]. The lower-price
volume comparison fails its pre-trend test.

\begin{center}
\includegraphics[page=4,width=0.9\textwidth]{../input/main_spec_2026-09-22.pdf}

\includegraphics[page=5,width=0.9\textwidth]{../input/main_spec_2026-09-22.pdf}
\end{center}

Pre-2013 gentrification explains part, but not all, of the pattern. The new
\texttt{tract\_demographic\_change} audit measures each sale's 2010 tract change
from the 2000 Census to the 2008--2012 ACS. The measures are the bachelor's degree
share, non-Hispanic white and Black shares, and log real mean household income. In the top price quartile, closed-school areas had
gained more college graduates (21.0 against 18.9 percentage points) and
more white residents (14.7 against 10.3). They lost fewer Black
residents ($+2.0$ against $-10.4$). Year effects that vary with these four changes
leave the slope almost unchanged: 0.139 [$-0.01$, 0.28] for all property and 0.114
[$-0.02$, 0.25] for single-family. They halve the top-quartile effect, to 0.102 [$-0.07$,
0.27], and remove it for single-family homes ($-0.014$). No single
neighborhood drives the gradient. Dropping each of the seven top-quartile closed
sites (Stewart, Trumbull, Peabody, Near North, Lafayette, King, Duprey/Humboldt)
leaves the hedonic slope between 0.133 and 0.227. Dropping Trumbull raises it.

The defensible statement is therefore limited. After 2013, closed-school neighborhoods did
relatively better the more expensive they were beforehand. Part of that coincides
with gentrification measured before the closures. The remaining gradient could reflect the closures, including quick reuse of
vacated buildings in hot markets, or gentrification that the four tract measures
miss. These data cannot yet separate the two. Coding each closed building's disposition and timing is the
next step for that question. The intervals are wide throughout; the absence of
an average effect is not a precise zero.

\small
Reproduction: \texttt{make} in \texttt{tasks/audits/home\_prices\_twfe/code}
(models and the robustness packet \texttt{event\_studies.pdf}; figures here from
\texttt{main\_spec.pdf}, frozen in \texttt{logbook/exhibits/home\_prices\_twfe\_2026-09-22/}),
\texttt{tasks/audits/home\_sales\_volume/code}, and
\texttt{tasks/audits/tract\_demographic\_change/code}. Census vintages: 2000 SF3
and 2008--2012 ACS via the Census API, retrieved September 22, 2026.
\normalsize
